\documentclass[11pt]{article}
\usepackage[margin=1in]{geometry}
\usepackage{amsmath,amssymb,amsthm}
\usepackage{graphicx}
\usepackage{microtype}
\usepackage{booktabs}
\usepackage{titlesec}
\usepackage[hidelinks]{hyperref}
\titleformat{\section}{\normalfont\Large\bfseries}{\thesection}{0.8em}{}
\titleformat{\subsection}{\normalfont\large\bfseries}{\thesubsection}{0.8em}{}

\newtheorem{theorem}{Theorem}
\newtheorem{proposition}[theorem]{Proposition}
\newtheorem{corollary}[theorem]{Corollary}
\newtheorem{remark}[theorem]{Remark}
\newcommand{\R}{\mathbb{R}}
\newcommand{\as}{\operatorname{as}}
\newcommand{\IS}{\operatorname{IS}}
\newcommand{\OS}{\operatorname{OS}}
\newcommand{\vrad}{\operatorname{vrad}}
\newcommand{\Cin}{\mathcal C_{\mathrm{in}}}
\newcommand{\Cout}{\mathcal C_{\mathrm{out}}}

\title{Envelope--Santal\'o extension for extremal general affine surface area}
\author{Candidate full result for arXiv:2103.00294}
\date{19 July 2026}

\begin{document}
\maketitle

\noindent\textbf{Status.} Candidate full result, likely valid, pending
specialist review.  The intended concave-side program is completed for a new
nonmonotone class.  Two printed claims in the source require correction, so
the theorem below proves the outer statements directly rather than using the
source's polarity argument.

\section{Problem and outcome}

Hoehner studies extremal general affine surface areas under global
monotonicity of $\varphi(t)/\sqrt t$ and asks for the results of the article
for a richer class than
$\operatorname{Conc}^{-}(0,\infty)\cup
\operatorname{Conc}^{+}(0,\infty)$ \cite{Hoehner}.  The exact source prompt is
reproduced below.

\begin{quote}
``Ultimately, it is an open problem to obtain the results of this article for
a richer class of functions than
$\operatorname{Conc}^{-}(0,\infty)\cup\operatorname{Conc}^{+}(0,\infty)$.''
\end{quote}

\begin{center}
\includegraphics[width=\textwidth]{figures/open_problem_crop.png}
\end{center}

We introduce an \emph{interior-peaked} class, prove existence, attainment,
Hausdorff continuity, sharp affine-isoperimetric envelopes, and both inner and
outer Blaschke--Santal\'o inequalities, and determine the exact equality
conditions.  The class contains
\[
  \varphi(t)=\log(1+t),
\]
whose quotient by $\sqrt t$ is nonmonotone, so it is genuinely outside both
source classes.  Equality is necessarily different from the source theory:
this example has many non-ellipsoidal equality bodies.

The proof also exposes two independent issues in the published statements.
First, the claimed ``equality iff ellipsoid'' assertion in the source
Blaschke--Santal\'o theorem fails on a scaled ball for the paper's own
arctangent example.  Second, centroid zero is not preserved by polarity, so
the asserted bijection between centered inner and centered outer competitors
is false.  Neither issue affects the direct proofs below.

\section{Definitions and source inputs}

Let $B_n$ be the Euclidean unit ball in $\R^n$ and let
\[
 \vrad(K)=\left(\frac{|K|}{|B_n|}\right)^{1/n}.
\]
All ambient bodies contain the origin in their interiors.  Write
\[
 \Cin(K)=\{C\subset K:C\text{ is a convex body and }g(C)=0\},
\]
\[
 \Cout(K)=\{D\supset K:D\text{ is a convex body and }g(D)=0\}.
\]
For $\varphi\in\operatorname{Conc}(0,\infty)$, put
\[
 \IS_\varphi(K)=\sup_{C\in\Cin(K)}\as_\varphi(C).
\]
Define $\varphi^*(t)=t\varphi(1/t)$ and retain the source's original outer
functional
\[
 \OS_{\varphi^*}(K)=\sup_{D\in\Cout(K)}\as_{\varphi^*}(D).
\]
The involution preserves $\operatorname{Conc}(0,\infty)$.

We use three standard facts from Ludwig's general affine surface-area theory
\cite{Ludwig}.  For a centered convex body $C$ and
$\theta\in\operatorname{Conc}(0,\infty)$,
\begin{equation}\label{eq:ludwig}
 \as_\theta(C)\leq\as_\theta(\vrad(C)B_n),
\end{equation}
with equality for strictly increasing $\theta$ only when $C$ is an
ellipsoid; $\as_\theta$ is upper semicontinuous under Hausdorff convergence;
and
\begin{equation}\label{eq:ball}
 \as_\theta(rB_n)=n|B_n|\frac{\theta(r^{-2n})}{\sqrt{r^{-2n}}}.
\end{equation}

Set
\[
 h_\varphi(t)=\frac{\varphi(t)}{\sqrt t},\qquad
 L_\varphi=\limsup_{t\to\infty}h_\varphi(t),\qquad
 H_\varphi(a)=\sup_{t\geq a}h_\varphi(t).
\]
We call $\varphi$ \emph{tail-gapped} if
\begin{equation}\label{eq:tailgap}
 L_\varphi<\infty
 \quad\text{and}\quad
 H_\varphi(a)>L_\varphi\quad\text{for every }a>0.
\end{equation}
It is \emph{interior-peaked} if it is tail-gapped and
\begin{equation}\label{eq:peak}
 M_\varphi:=\sup_{t>0}h_\varphi(t)=H_\varphi(1).
\end{equation}
Thus a global radial maximum is accessible at a finite parameter $t\geq1$.
Indeed, $H_\varphi(1)>L_\varphi$ implies that the supremum is attained.  Let
\begin{equation}\label{eq:argmax}
 \mathcal A_\varphi=\{t>0:h_\varphi(t)=M_\varphi\}.
\end{equation}

The basic identity for the dual quotient is
\begin{equation}\label{eq:dualh}
 \frac{\varphi^*(t)}{\sqrt t}
 =\sqrt t\,\varphi(1/t)=h_\varphi(1/t).
\end{equation}

\section{\hspace{0.35em}Idea of the proof}

Affine isoperimetry reduces both extremal problems to the same
one-dimensional tail envelope.  An inner competitor of volume radius $r$
uses the parameter $r^{-2n}$, whereas an outer $\varphi^*$ competitor uses
$r^{2n}$ by \eqref{eq:dualh}.  Finite tail limsup bounds both suprema.  The
strict tail gap is the compactness mechanism: it places a finite-radius ball
strictly above the value available to an inner sequence collapsing to zero
volume or an outer sequence escaping to infinite volume.  Blaschke selection
then gives both maximizers.

For continuity, concavity supplies a two-sided estimate under dilation; this
is stronger than the one-sided estimate used in the source.  Finally, the
interior-peak condition says that every radial envelope is bounded by the
value already accessible at the unit ball.  It therefore proves both
Blaschke--Santal\'o products factor by factor, without the source's invalid
polar-centroid identification.  Equality can then be read off from equality
in Ludwig's affine isoperimetric inequality and membership in the maximizing
set $\mathcal A_\varphi$.

\section{Two-sided dilation control}

The next elementary estimate is useful because it gives continuity for both
supremal functionals without polarizing their competitors.

\begin{proposition}\label{prop:dilation}
If $\theta\in\operatorname{Conc}(0,\infty)$, $D$ is a convex body containing
the origin, and $c\geq1$, then
\begin{equation}\label{eq:dilation}
 c^{-n}\as_\theta(D)\leq\as_\theta(cD)\leq c^n\as_\theta(D).
\end{equation}
Consequently, whenever the four quantities are finite,
\begin{align}
 c^{-n}\IS_\theta(K)&\leq\IS_\theta(cK)\leq c^n\IS_\theta(K),
 \label{eq:dilationIS}\\
 c^{-n}\OS_\theta(K)&\leq\OS_\theta(cK)\leq c^n\OS_\theta(K).
 \label{eq:dilationOS}
\end{align}
\end{proposition}

\begin{proof}
Every positive concave function on $(0,\infty)$ with value $0$ at the origin
is nondecreasing; otherwise a negative secant slope and concavity would force
the function eventually to become negative.  The change of scale in the
defining integral gives
\[
 \as_\theta(cD)
 =c^n\int_{\partial D}\theta(c^{-2n}q_D(x))
   \langle x,N_D(x)\rangle\,d\mu_D(x),
\]
where $q_D$ is the centro-affine curvature ratio.  For
$\lambda=c^{-2n}\in(0,1]$, monotonicity and concavity at the origin give
\[
 \lambda\theta(u)\leq\theta(\lambda u)\leq\theta(u).
\]
This proves \eqref{eq:dilation}.  The maps $C\mapsto cC$ and $D\mapsto cD$
are bijections between the corresponding inner and outer competitor
families, so taking suprema proves \eqref{eq:dilationIS} and
\eqref{eq:dilationOS}.
\end{proof}

\section{Full radial-envelope theorem}

\begin{theorem}\label{thm:envelopes}
Let $\varphi\in\operatorname{Conc}(0,\infty)$ be tail-gapped.  For every
convex body $K$ containing the origin in its interior:
\begin{enumerate}
 \item $\IS_\varphi(K)$ and $\OS_{\varphi^*}(K)$ are positive and finite,
 are attained by centered convex bodies, and depend continuously on $K$ in
 the Hausdorff metric;
 \item the sharp envelope inequalities hold:
 \begin{align}
  \IS_\varphi(K)
  &\leq n|B_n|H_\varphi(\vrad(K)^{-2n}),
  \label{eq:innerenv}\\
  \OS_{\varphi^*}(K)
  &\leq n|B_n|H_\varphi(\vrad(K)^{2n});
  \label{eq:outerenv}
 \end{align}
 \item equality holds in \eqref{eq:innerenv} and \eqref{eq:outerenv} for
 every centered Euclidean ball $K=RB_n$.
\end{enumerate}
\end{theorem}

\begin{proof}
The inner assertion is included for completeness.  If $C\in\Cin(K)$ and
$s_C=\vrad(C)^{-2n}$, then $s_C\geq\vrad(K)^{-2n}$.  Equations
\eqref{eq:ludwig} and \eqref{eq:ball} give
\[
 \as_\varphi(C)\leq n|B_n|h_\varphi(s_C)
 \leq n|B_n|H_\varphi(\vrad(K)^{-2n}).
\]
This proves finiteness and \eqref{eq:innerenv}.  If a maximizing sequence
collapsed to a lower-dimensional set, then $s_C\to\infty$ and its limiting
value would be at most $n|B_n|L_\varphi$.  A centered ball $\rho B_n\subset K$
and the strict tail gap supply a smaller centered ball in $K$ with value
strictly larger than $n|B_n|L_\varphi$.  Collapse is impossible.  Blaschke
selection, centroid continuity, and upper semicontinuity then give an inner
maximizer.

For the outer assertion, let $D\in\Cout(K)$ and put
$u_D=\vrad(D)^{2n}$.  By \eqref{eq:dualh},
\[
 \as_{\varphi^*}(D)
 \leq n|B_n|h_\varphi(u_D)
 \leq n|B_n|H_\varphi(\vrad(K)^{2n}),
\]
which proves finiteness and \eqref{eq:outerenv}.

Let $S=\OS_{\varphi^*}(K)$.  Choose $R$ with $K\subset RB_n$.  The strict
gap $H_\varphi(R^{2n})>L_\varphi$ supplies $u\geq R^{2n}$ with
$h_\varphi(u)>L_\varphi$.  Hence the centered ball
$u^{1/(2n)}B_n\supset K$ shows
\begin{equation}\label{eq:outergap}
 S>n|B_n|L_\varphi.
\end{equation}
Let $D_j$ be an outer maximizing sequence.  Its volumes are bounded: an
unbounded-volume subsequence would have $u_{D_j}\to\infty$ and therefore,
by the preceding affine-isoperimetric estimate, limsup at most
$n|B_n|L_\varphi$, contradicting \eqref{eq:outergap}.  Fix
$\rho B_n\subset K$.  Since every $D_j$ contains $\rho B_n$, the cone with
apex $x\in D_j$ and base $\rho B_n\cap x^\perp$ lies in $D_j$, so
\[
 |D_j|\geq \frac{|x|\rho^{n-1}|B_{n-1}|}{n}.
\]
The volume bound therefore gives a uniform diameter bound.  Blaschke
selection produces a Hausdorff limit $D_0\supset K$; it is full-dimensional,
has centroid zero, and upper semicontinuity gives
$S\leq\as_{\varphi^*}(D_0)\leq S$.  Thus the outer supremum is attained.

For $K=RB_n$, every parameter $s\geq R^{-2n}$ is realized by the inner ball
$s^{-1/(2n)}B_n$, and every $u\geq R^{2n}$ is realized by the outer ball
$u^{1/(2n)}B_n$.  This proves both sharpness statements.

It remains to prove continuity.  Both functionals are monotone under
inclusion, with $\IS$ increasing and $\OS$ decreasing.  If $K_j\to K$, then
there are $c_j,d_j\downarrow1$ such that
$K_j\subset c_jK$ and $K\subset d_jK_j$.  Inclusion monotonicity and
Proposition \ref{prop:dilation} give
 \[
 d_j^{-n}\IS_\varphi(K)\leq\IS_\varphi(K_j)
 \leq c_j^n\IS_\varphi(K)
 \]
and
\[
 c_j^{-n}\OS_{\varphi^*}(K)\leq\OS_{\varphi^*}(K_j)
 \leq d_j^n\OS_{\varphi^*}(K).
\]
Both squeeze limits prove Hausdorff continuity.
\end{proof}

\section{Blaschke--Santal\'o theorem and exact equality}

\begin{theorem}\label{thm:full}
Let $\varphi$ be interior-peaked.  Then for every convex body $K$ containing
the origin in its interior,
\begin{align}
 \IS_\varphi(K)&\leq\IS_\varphi(B_n)=n|B_n|M_\varphi,
 \label{eq:univinner}\\
 \OS_{\varphi^*}(K)&\leq\OS_{\varphi^*}(B_n)=n|B_n|M_\varphi.
 \label{eq:univouter}
\end{align}
Consequently,
\begin{align}
 \IS_\varphi(K)\IS_\varphi(K^\circ)&\leq\IS_\varphi(B_n)^2,
 \label{eq:BSinner}\\
 \OS_{\varphi^*}(K)\OS_{\varphi^*}(K^\circ)
 &\leq\OS_{\varphi^*}(B_n)^2.
 \label{eq:BSouter}
\end{align}
No centroid assumption on $K$ is needed.

If $\varphi$ is strictly increasing, equality in \eqref{eq:univinner} holds
if and only if $K$ contains a centered ellipsoid $E$ such that
\begin{equation}\label{eq:innereq}
 \vrad(E)^{-2n}\in\mathcal A_\varphi.
\end{equation}
If $\varphi^*$ is strictly increasing, equality in \eqref{eq:univouter}
holds if and only if $K$ is contained in a centered ellipsoid $D$ such that
\begin{equation}\label{eq:outereq}
 \vrad(D)^{2n}\in\mathcal A_\varphi.
\end{equation}
Thus equality in either product inequality holds exactly when the
corresponding containment condition holds for both $K$ and $K^\circ$.
\end{theorem}

\begin{proof}
By definition of the interior-peaked class,
$H_\varphi(a)\leq M_\varphi=H_\varphi(1)$ for every $a>0$.  Theorem
\ref{thm:envelopes}, applied to $K$ and to $B_n$, proves
\eqref{eq:univinner} and \eqref{eq:univouter}.  Applying each universal bound
to $K$ and $K^\circ$ proves the two product inequalities.

Suppose $\varphi$ is strictly increasing and equality holds in
\eqref{eq:univinner}.  Let $C\in\Cin(K)$ be an inner maximizer from Theorem
\ref{thm:envelopes}.  The chain
\[
 n|B_n|M_\varphi=\as_\varphi(C)
 \leq n|B_n|h_\varphi(\vrad(C)^{-2n})
 \leq n|B_n|M_\varphi
\]
is equality throughout.  Ludwig's equality condition makes $C$ an ellipsoid,
and its parameter belongs to $\mathcal A_\varphi$.  The converse follows by
using such an ellipsoid as a competitor.  The outer proof is identical with
an outer maximizer, \eqref{eq:dualh}, and strict increase of $\varphi^*$.
Finally, two positive numbers each bounded by the same positive constant have
maximal product exactly when both attain that constant.  This gives the
product equality criteria.
\end{proof}

\section{The logarithmic class member and non-rigid equality}

\begin{corollary}\label{cor:log}
Let $\varphi(t)=\log(1+t)$.  There is a unique $t_*>1$ satisfying
\begin{equation}\label{eq:tstar}
 \log(1+t_*)=\frac{2t_*}{1+t_*},
\end{equation}
and $t_*=3.921553634568\ldots$.  The function $\varphi$ is
interior-peaked, belongs to neither source monotonicity class, and
$\mathcal A_\varphi=\{t_*\}$.  Both $\varphi$ and $\varphi^*$ are strictly
increasing, so all conclusions and equality conditions of Theorem
\ref{thm:full} apply.

Put $r_*=t_*^{-1/(2n)}<1$.  Then
\begin{align*}
 \IS_\varphi(K)=\IS_\varphi(B_n)
 &\Longleftrightarrow K\text{ contains a centered ellipsoid of volume radius }r_*,\\
 \OS_{\varphi^*}(K)=\OS_{\varphi^*}(B_n)
 &\Longleftrightarrow K\text{ is contained in a centered ellipsoid of volume radius }r_*^{-1}.
\end{align*}
For $n\geq2$ and
\begin{equation}\label{eq:epsilon}
 0<\varepsilon<r_*^{-1}-1,
\end{equation}
the capsule
\[
 K_\varepsilon=B_n+\varepsilon[-e_1,e_1]
\]
is not an ellipsoid, yet equality holds for $K_\varepsilon$ in both product
inequalities \eqref{eq:BSinner} and \eqref{eq:BSouter}.
\end{corollary}

\begin{proof}
Differentiation shows that the sign of
$h_\varphi'(t)$ is the sign of
\[
 q(t)=\frac{2t}{1+t}-\log(1+t).
\]
Since $q'(t)=(1-t)/(1+t)^2$, the function $q$ increases on $(0,1)$,
decreases on $(1,\infty)$, and tends to $-\infty$.  It has the unique
positive zero $t_*>1$, so $h_\varphi$ increases up to $t_*$ and decreases
afterward.  Moreover, $h_\varphi(t)\to0$ at both endpoints.  Hence
$\varphi$ is interior-peaked, globally nonmonotone, and has the stated unique
maximizer.  The function $\varphi$ is strictly increasing, while
\[
 (\varphi^*)'(t)=\log(1+1/t)-\frac1{1+t}>0
\]
because $\log(1+x)>x/(1+x)$ for $x>0$.

The equality characterizations now follow from Theorem \ref{thm:full}.  For
the explicit capsule,
\[
 r_*B_n\subset B_n\subset K_\varepsilon
 \subset(1+\varepsilon)B_n\subset r_*^{-1}B_n.
\]
Polarity reverses these inclusions, so the same two-sided containment holds
for $K_\varepsilon^\circ$.  Thus both $K_\varepsilon$ and its polar satisfy
both equality criteria.  The capsule has nontrivial cylindrical boundary
segments and is not an ellipsoid.
\end{proof}

\section{Two corrections to the source paper}

\subsection{The printed ellipsoid equality claim is false}

The source's Theorem 6.4 asserts, under strict increase of $\varphi$, that
equality in its Blaschke--Santal\'o inequality holds if and only if $K$ is an
ellipsoid.  Its own family
\[
 \varphi_m(t)=\arctan(t^{1/m}),\qquad m\geq2,
\]
contradicts the ``if'' direction.  The source proves that $h_{\varphi_m}$ is
strictly decreasing and that
\begin{equation}\label{eq:atanstrict}
 \varphi_m(s)\varphi_m(1/s)<\varphi_m(1)^2
 \qquad(s\neq1).
\end{equation}
Indeed, if $a=\arctan(s^{1/m})$, then the left side is
$a(\pi/2-a)$, whose unique maximum occurs at $a=\pi/4$, equivalently $s=1$.

Choose the ellipsoid $K=rB_n$ with $r\neq1$ and set $s=r^{-2n}$.  Since
$h_{\varphi_m}$ is decreasing, the sharp ball formula gives
\[
 \IS_{\varphi_m}(rB_n)\IS_{\varphi_m}(r^{-1}B_n)
 =(n|B_n|)^2\varphi_m(s)\varphi_m(1/s)
 <(n|B_n|)^2\varphi_m(1)^2
 =\IS_{\varphi_m}(B_n)^2.
\]
Thus a scaled Euclidean ball is an ellipsoid but is not an equality case.
The exact envelope equality criterion in Theorem \ref{thm:full} avoids this
homogeneity error.

\subsection{Centered competitor families are not polar duals}

The source states that $C\in\Cin(K)$ if and only if
$C^\circ\in\Cout(K^\circ)$.  Under the printed definitions this would require
$g(C)=0$ if and only if $g(C^\circ)=0$, which is false.

An exact planar counterexample is the quadrilateral
\[
 C=\operatorname{conv}\left\{
 \frac{(-7,-4)}9,\frac{(11,-4)}9,
 \frac{(2,5)}9,\frac{(-7,5)}9\right\}.
\]
The polygon-centroid formula gives $g(C)=(0,0)$.  Directly intersecting the
four polar half-planes gives
\[
 C^\circ=\operatorname{conv}\left\{
 (0,-9/4),(9/7,9/7),(0,9/5),(-9/7,0)\right\},
\]
whose centroid is
\[
 g(C^\circ)=(0,9/140)\neq0.
\]
\par\noindent
Taking $K=C$ disproves the asserted competitor-family equivalence and hence
the extremal polarity identity derived from it.  This is why Theorems
\ref{thm:envelopes} and \ref{thm:full} treat the original outer functional
directly.

\section{Scope, verification, and novelty}

Theorem \ref{thm:full} supplies the complete intended chain---existence,
attainment, continuity, sharp affine bounds, inner and outer
Blaschke--Santal\'o inequalities, and exact equality---for a concrete third
class outside the source monotonicity union.  It does not assert that every
tail-gapped function satisfies the normalized product inequality; the global
peak condition is precisely what makes both original functionals uniformly
bounded by their unit-ball values.

The proof was checked against collapse and escape of maximizing sequences,
centroid preservation under Hausdorff limits, both directions of the dilation
estimate, and equality in every inequality of the maximizing chain.  The
script \texttt{code/check\_examples.py} verifies the reported logarithmic
root, strictness of the arctangent scaled-ball example, and both rational
polygon centroids.  The computations are not used in place of analytic
proofs.

A bounded novelty search covered the four local run indexes, exact title and
phrase searches, the arXiv record, the journal version, and works listed as
citing the paper through 19 July 2026.  No erratum, later solution of the open
problem, radial-envelope theorem, or published notice of the two corrections
was found.  Absolute novelty is not certified.

\paragraph{Human review recommendation.}
Send to a convex-geometric analyst.  The highest-value checks are the direct
outer compactness proof, the two-sided dilation identity for nonsmooth bodies,
and the interpretation of ``full'' as the complete intended concave-side
program for a genuinely new class rather than a verbatim preservation of the
source's two incorrect statements.

\begin{thebibliography}{9}
\bibitem{Hoehner}
S. Hoehner,
\emph{Extremal general affine surface areas},
J. Math. Anal. Appl. 505 (2022), 125506;
\href{https://arxiv.org/abs/2103.00294}{arXiv:2103.00294}.

\par\medskip

\bibitem{Ludwig}
M. Ludwig,
\emph{General affine surface areas},
Adv. Math. 224 (2010), 2346--2360.
\end{thebibliography}

\end{document}
